\documentclass[11pt]{article}
\usepackage[letterpaper,margin=2cm]{geometry}
\usepackage[T1]{fontenc}
\usepackage{lmodern}
\usepackage{microtype}
\usepackage{booktabs}
\usepackage[compact]{titlesec}
\usepackage{enumitem}
\usepackage{xcolor}
\usepackage{graphicx}
\usepackage[colorlinks,urlcolor=blue!60!black,linkcolor=black]{hyperref}

\titlespacing*{\section}{0pt}{1.6ex}{0.6ex}
\setlist{nosep,leftmargin=*}
\setlength{\parindent}{0pt}
\setlength{\parskip}{0.6ex}

% Branch the links point to.
\newcommand{\branch}{main}
\newcommand{\repo}{https://github.com/Shivsi7/dynamic-grasping/blob/\branch}
% Where the example videos live.
\newcommand{\videos}{\repo/docs/report/videos}

% Strip of four frames for one example (frames/<scene>_0..3.jpg): the scene at t = 0, then (A), (B), (C).
\newcommand{\fr}[1]{\includegraphics[width=0.215\textwidth]{frames/#1}}
\newcommand{\strip}[1]{\par\noindent\fr{#1_0}\hfill\fr{#1_1}\hfill\fr{#1_2}\hfill\fr{#1_3}\par\medskip}

\begin{document}

\begin{center}
  {\Large\bfseries Dynamic grasping: progress report}\\[0.6ex]
  [Authors] \quad$\cdot$\quad 6.4212, Fall 2026 \quad$\cdot$\quad \today
\end{center}

\section{What we did}

\textbf{Environment.} In Drake, a KUKA iiwa with a WSG gripper picks a YCB object off a cart moving at
0.05--0.3\,m/s, while 0--4 spheres and boxes cross the workspace, on wheeled rigs or lobbed. Scenes come
from a seeded generator (always solvable by getting out of the way) and from preset grids that block the
direct grasp (a post, a plank, a pillar).

\textbf{Evaluation.} Compute time costs simulated time (measured latency). Each episode runs to the end
of the scene: \emph{success} (grasped, lifted, still held at the end), \emph{collision}, \emph{dropped}
or \emph{timeout}. 214 scenes: dev 20, preset grids 114 (48 + 48 + 18), held-out 80 (a separate random
seed; some of its failures were inspected during development).

\textbf{(A) Reactive QP.} Every 10\,ms, a QP over joint velocities tracks a velocity-matched grasp pose,
with velocity dampers against the obstacles' predicted swept volume over the next 0.2\,s. Fast but local.
\href{\repo/docs/planner_A.md}{Details}.

\textbf{(B) Space-time RRT-Connect.} Plans in $(q,t)$ to a pre-grasp at an intercept time, then a
fixed grasp-and-lift motion. Replans every 0.5\,s, within acceleration and torque limits, with escape
moves when the plan turns unsafe. \href{\repo/docs/planner_B.md}{Details}.

\textbf{(C) Space-time trajectory optimization.} Plans once, before the object is in reach. IK picks a
collision-free grasp; Drake's kinematic trajectory optimization (B-spline, free duration) plans the
approach, clear of the obstacles' predicted motion. A second solve closes, lifts and holds until the scene
ends. \href{\repo/docs/planner_C.md}{Details}.

\section{Differences from the plan}

\begin{itemize}
  \item \textbf{Environment:} besides the planned tall post, added plank and pillar grids that block the
    direct grasp.
  \item \textbf{Evaluation:} compute time costs simulated time, and a success must hold the object to
    the end of the scene; neither was planned. Fixed scene sets, with a held-out set.
  \item \textbf{(A):} as planned; the swept volume is sampled at $t$, $t+0.1$, $t+0.2$\,s instead of
    using capsules.
  \item \textbf{(B):} written from scratch on Drake's distance queries, without an existing planner or
    VAMP. Added acceleration and torque limits, several grasp candidates, escape moves, and collision
    checks of the held object.
  \item \textbf{(C):} an optional extension, now a full method. Seeded by IK, not by (A). Plans once
    instead of re-solving at a few Hz, and also plans the hold.
\end{itemize}

\section{Results}

\begin{center}
\begin{tabular}{@{}lccc@{}}
  \toprule
  & (A) reactive & (B) RRT & (C) traj.\ opt. \\
  \midrule
  Success, $n$ (\%)                                           & 132 (62\%) & 173 (81\%) & 169 (79\%) \\
  Failures: collision / timeout / dropped                     & 6 / 71 / 5 & 0 / 41 / 0 & 8 / 36 / 1 \\
  Time to grasp after the object comes into reach (median, s) & 2.21       & 2.66       & 1.38 \\
  Min.\ obstacle clearance (median, cm)                       & 5.3        & 5.5        & 4.0 \\
  Compute per decision (median, ms)                           & 2.2        & 10         & 841 \\
  Peak torque (median, fraction of limit)                     & 2.20       & 0.34       & 0.43 \\
  Peak joint speed (median, fraction of limit)                & 0.60       & 0.44       & 0.57 \\
  \bottomrule
\end{tabular}

\smallskip
{\small All 214 scenes, measured latency, scored to the end of each scene; held-out 80 from a separate seed.}
\end{center}

\clearpage
\section{Examples}

Four held-out scenes; frames from each video: the scene at $t=0$, then (A), (B), (C).

\textbf{All three succeed.} \texttt{mix\_0074}: all three grasp the cracker box among four moving obstacles,
1.0 to 1.7\,s after it comes into reach, (C) first. \href{\videos/mix_0074.mp4}{Video}.
\strip{mix_0074}

\textbf{All three fail.} \texttt{mix\_0079}: three boxes and a sphere crowd the arm; (A) cannot get past
them, (B) escapes and times out, and (C)'s plans arrive too late, so a box hits the parked arm. \href{\videos/mix_0079.mp4}{Video}.
\strip{mix_0079}

\textbf{Only (A) fails.} \texttt{mix\_0036}: all three grasp the mustard bottle, but it slips out of
(A)'s gripper right after the lift. \href{\videos/mix_0036.mp4}{Video}.
\strip{mix_0036}

\textbf{Only (C) succeeds.} \texttt{mix\_0043}: two boxes and a sphere pass the arm; (B) never leaves its
escape pose, (A) grasps late and flings the bottle away swerving around the sphere, and (C) holds it. \href{\videos/mix_0043.mp4}{Video}.
\strip{mix_0043}

\section{Next steps}

\begin{enumerate}
  \item More variability in the scenes, e.g.\ objects flying around.
  \item Keep optimizing the methods, mainly (A) and (B); a lot of time already went into (C).
  \item Reinforcement learning as a fourth method.
  \item Perception, for a more realistic setting: the hardest, and the lowest priority.
\end{enumerate}

\end{document}
